
% Esta seção trata a construção do documento, especificamente a disposição dos capítulos e uma introdução do que está sendo abordado nestes.

\section{Apresentação do Trabalho}
    \label{sec:apresentacao-trabalho}

    No capítulo seguinte, fundamentação teórica, são discutidos assuntos relacionados ao trabalho, como Análise e visualização de dados, Aprendizado de Máquina, Engenharia de Software e as linguagens e bibliotecas utilizadas no desenvolvimento do sistema proposto.

    No terceiro capítulo são apresentados tópicos relacionados a modelagem do sistema, contendo os requisitos, funcionais e não funcionais, diagramas, modelo de dados e a arquitetura modular do sistema.

    O quarto capítulo é dedicado a implementação do sistema, contendo a descrição do processo de desenvolvimento, a arquitetura do sistema, a descrição das tecnologias utilizadas e a descrição do processo de desenvolvimento.

    Adiante, no quinto capítulo são apresentados os resultados obtidos com o sistema, contendo a descrição da interface, exemplos de uso e a avaliação do sistema.

    No sexto capítulo, são apresentadas a conclusão obtida com o presente trabalho, bem como sugestões de trabalhos futuros. Por fim, são listadas as referências bibliográficas utilizadas durante este trabalho, seguidas por um apêndice com matérias relacionados ao trabalho e um anexo com o código-fonte do sistema.
